\section{绪论}
\subsection{研究背景及意义}
随着全球能源危机的加剧与“双碳”目标的推进，发展清洁可再生能源已成为电力系统的核心战略。光储充电站结合了光伏发电（Photovoltaic, PV）、储能系统（Energy Storage System, ESS）与电动汽车充电设施，通过交直流微电网技术实现了源、网、荷、储的深度协同。相比于传统的交流配电网，直流微电网省去了中间的交直流转换环节，极大提高了分布式电源与直流负载的接入效率。光储充电站不仅能够实现光伏电能的就地消纳，有效平抑光照环境引发的功率输出随机波动，还能通过储能系统的“削峰填谷”显著降低大功率充电负荷对配电网的冲击\upcite{chen202404025,1020040865nh}。 


然而，光储充直流微电网本质上是一个由海量非线性电力电子变换器构成的强耦合复杂系统，其安全稳定运行面临着多源大扰动的严峻挑战。在光伏侧，光伏发电受气象条件制约，极易发生光伏辐照度降低导致的输出功率骤降；在充电桩负载侧，电动汽车充电桩在高压快速充电工况下，其内部高带宽的闭环控制使其表现出典型的恒功率负载（Constant Power Load, CPL）特性。CPL具有独特的负增量阻抗（Negative Incremental Impedance）效应，即当直流母线电压跌落时，为维持恒定功率输出，负载会自动增大电流抽取。当“源”侧的光伏出力骤降与“荷”侧的大功率突加叠加时，CPL的负阻抗正反馈机制将被快速放大，容易引发直流母线电压的振荡从而导致系统崩溃。这就要求“储”侧的蓄电池系统及其双向变换器必须具备良好的暂态响应与功率支撑能力，以维持系统能量平衡。


为了深入探究光储充电站系统在上述复杂多源扰动下的暂态机理，本文建立了涵盖光伏、储能、网侧逆变器及等效恒功率负载的系统级暂态分析模型。在科学、合理的系统建模基础上，全面剖析光伏辐照度突降、大功率负载切入与切出、三相电网短路等多重极端工况下，各变换器之间的暂态交互机理与母线电压演变波形。本课题的研究成果能够在为光储充电站的容量配置与系统规划提供理论支撑的同时，也能为保障微电网的安全稳定运行奠定坚实的工程仿真基础。
\subsection{国内外研究现状}

随着光储充直流微电网系统的快速发展，其内部因接入海量非线性电力电子设备而引发的系统级稳定性问题，引起了国内外学者的广泛关注。特别是在实际复杂工况下，源侧光伏发电极易受极端气象条件制约而发生功率骤降，荷侧大功率快充设施表现出严密的恒功率负阻抗特性，加之网侧可能遭遇的大电网电压跌落或低压穿越（LVRT）等暂态故障，这些多源异构扰动的交织与耦合，使得微电网的瞬时功率平衡与直流母线电压极易遭到破坏。目前，针对光储充直流微电网系统的暂态稳定性研究，主要集中在系统等效降阶建模、大扰动非线性稳定性分析以及暂态主动支撑控制策略三个核心方面。

\textbf{1. 直流微电网与恒功率负载的等效建模}

在系统级仿真与分析中，直接采用包含高频开关管的物理拓扑（如全桥LLC变换器）会导致严重的维度灾难与计算资源枯竭。因此，国内外文献普遍采用平均值模型或大信号等效模型对系统进行降阶处理。对于电动汽车充电桩等闭环受控变换器，研究表明在其控制带宽内，负载端会表现出严密的恒功率特性。在研究直流微电网中恒功率负荷的动态特性时，剥离开关纹波，利用受控电流源对CPL进行宏观外特性等效，是进行大尺度暂态分析的高效手段\upcite{chen2019cpl}。

\textbf{2. CPL负阻抗特性与大扰动暂态分析}

恒功率负载由于具有负增量阻抗特性，会显著削弱直流微电网的系统阻尼。当系统遭遇光照突降或大功率负载突加等大扰动时，CPL的负阻抗正反馈机制极易引发母线电压的剧烈振荡甚至系统崩溃\upcite{sun2021cplstability}。为了精确评估系统在大扰动下的稳定边界，国内外学者通常脱离传统的小信号线性化方法，转而采用大信号分析法对系统进行暂态分析。

\textbf{3. 提升暂态稳定性的控制策略与经济调度}

为了抑制微电网的波动从而降低失稳风险并提升运行经济性，当前的控制策略主要分为底层的暂态有源阻尼与上层的稳态经济调度两类。优化光储充一体化直流微电网的系统架构与底层控制策略，能够有效增强变换器对母线电压的动态支撑能力\upcite{wang2020microgrid}；低电压穿越（LVRT）等主动控制策略因加强了并网逆变器对外部电网故障的应对能力而得到深入研究与广泛应用\upcite{yang2019lvrt}；随着系统规模扩大，改进粒子群算法（PSO）等启发式算法与光储充系统的能量管理策略相结合，通过充分计及电池寿命的衰减成本等上层经济调度手段，大幅降低系统运行成本。\upcite{zhang2021pso, li2022battery}。

\subsection{本文研究内容}

综合现有文献可知，尽管学术界对单一CPL的负阻抗特性及其抑制策略已有深入研究，但多数文献仅针对简化的“单源-单荷”直流系统。对于包含光伏阵列、储能系统、网侧逆变器及大功率充电桩的完整“光储充”多源微电网，在面对光伏出力骤降与负荷台阶突变等多重极端扰动叠加时，各变换器之间的暂态交互机理仍需进一步通过高效的系统级仿真进行验证与剖析。这也是本文开展暂态分析与日前经济调度的切入点。